\documentclass[11pt,a4paper]{article}
\usepackage[margin=1in]{geometry}
\usepackage{amsmath,amssymb,amsthm}
\usepackage{algorithm}
\usepackage{algpseudocode}
\usepackage{booktabs}
\usepackage{hyperref}

\theoremstyle{definition}
\newtheorem{definition}{\textbf{Definition}}[section]
\newtheorem{theorem}{\textbf{Theorem}}[section]
\newtheorem{lemma}[theorem]{\textbf{Lemma}}
\newtheorem{remark}{\textbf{Remark}}[section]

\title{\textbf{Reward Modeling and Bradley-Terry Preference Learning}}
\author{\textbf{Team Scaibu} \\ \small Email: \texttt{jaydeep.9664920749@gmail.com} \\ \small \textbf{Architecture and Core Engine Team}}
\date{\textbf{October 2026}}

\begin{document}
\maketitle

\section{Definitions and Cognitive Mental Models}

\subsection{Formal Mathematical Definition}
\textbf{Definition (Pairwise Preference Likelihood and Bradley-Terry Models):}
Reinforcement Learning from Human Feedback (RLHF) requires a scalar reward signal $r_\theta(x, y) \in \mathbb{R}$ that evaluates the alignment, helpfulness, and safety of a generated text response $y$ given a prompt $x$. Because asking humans to assign absolute numerical scores (e.g. ``$7.34$ out of $10$'') suffers from high intra-annotator and inter-annotator calibration variance, preference feedback is collected via comparative pairwise judgments: given prompt $x$ and candidate completions $(y_w, y_l)$, human judges declare that response $y_w$ is preferred over response $y_l$ ($y_w \succ y_l$).

Under the \textbf{Bradley-Terry Preference Model}, the probability that response $y_w$ is preferred over $y_l$ is modeled as a logistic sigmoid function of the latent reward difference:
{\boldmath
\begin{equation}
P(y_w \succ y_l \mid x) = \sigma\left( r_\theta(x, y_w) - r_\theta(x, y_l) \right) = \frac{1}{1 + \exp\left( -\left( r_\theta(x, y_w) - r_\theta(x, y_l) \right) \right)}
\end{equation}
}
The reward model parameters $\theta$ are trained by minimizing the negative log-likelihood over a dataset of pairwise comparisons $\mathcal{D} = \{(x, y_w, y_l)\}$:
{\boldmath
\begin{equation}
\mathcal{L}_{\text{RM}}(\theta) = - \mathbb{E}_{(x, y_w, y_l) \sim \mathcal{D}} \left[ \ln \sigma\left( r_\theta(x, y_w) - r_\theta(x, y_l) - m \right) \right]
\end{equation}
}
where $m \ge 0$ is an optional target preference margin enforcing separation. Pairwise accuracy is evaluated as $\mathbb{E}[\mathbb{I}(r_\theta(x, y_w) > r_\theta(x, y_l))]$.

\subsection{Expanded Non-Technical Definition (Intuition for Anyone)}
\textbf{Intuitive Conceptual Definition:}
\begin{enumerate}
    \item \textbf{Why Pairwise Comparisons Work}: If you ask two different judges to rate a coffee on a scale from 1 to 100, Judge A might give it 85 and Judge B might give it 60 because they have different internal scales. But if you give both judges two cups of coffee side-by-side and ask ``Which cup tastes better?'', both judges will agree almost 100\% of the time!
    \item \textbf{Converting Votes into Points}: Bradley-Terry is the exact mathematical model used to rank chess grandmasters (Elo rating). If Player W consistently beats Player L, the model assigns Player W a higher score so that their score difference predicts their win rate.
    \item \textbf{The Length Bias Trap}: A notorious pathology in reward modeling is length bias: longer responses often look more thorough, so untrained reward models start giving maximum points to 2,000-word rambling fluff. Modern reward modeling actively monitors and penalizes correlation with word count.
    \item \textbf{The RLHF Fuel}: Once trained, this scalar reward model becomes the scoring engine for PPO or rejection sampling, guiding the AI to generate helpful, truthful, and safe answers.
\end{enumerate}

\subsection{Child-Friendly Mental Model (Physical Metaphor)}
Imagine an Olympic figure skating judge:
\begin{itemize}
    \item \textbf{The Old Flawed Scorecard}: Trying to give a random score like ``9.82'' leads to arguments among judges.
    \item \textbf{The Side-by-Side Dual}: Skater A and Skater B perform back-to-back. The judge simply points: ``Skater A landed their jump cleanly, Skater B stumbled. Skater A wins!''
    \item \textbf{The Leaderboard}: Every time Skater A wins, their ranking points go up. Every time a skater loses, their points go down.
    \item \textbf{Result}: The leaderboard automatically reveals who the most skilled skaters are without anyone having to invent arbitrary numbers!
\end{itemize}

\section{Core Mathematical Equations and Definitions}

\subsection{Bradley-Terry Negative Log-Likelihood Loss}
{\boldmath
\begin{equation}
\mathcal{L}_{\text{RM}} = \frac{1}{B}\sum_{i=1}^B \ln\left( 1 + \exp\left( -(r_\theta(x_i, y_{w, i}) - r_\theta(x_i, y_{l, i}) - m) \right) \right)
\end{equation}
}
\begin{itemize}
    \item \textbf{Formal Definition}: The empirical mean softplus loss equivalent to $-\ln \sigma(\Delta r - m)$.
    \item \textbf{What It Means}: Drives the score of preferred responses higher than dispreferred responses with smooth, non-saturating gradients.
    \item \textbf{Why It Is Important}: Primary optimization objective for training language model reward estimators in alignment pipelines.
\end{itemize}

\subsection{Pairwise Classification Accuracy}
{\boldmath
\begin{equation}
\operatorname{Acc}_{\text{pair}} = \frac{1}{B}\sum_{i=1}^B \mathbb{I}\left( r_\theta(x_i, y_{w, i}) > r_\theta(x_i, y_{l, i}) \right)
\end{equation}
}
\begin{itemize}
    \item \textbf{Formal Definition}: Fraction of evaluation pairs where the reward model correctly predicts human preference.
    \item \textbf{What It Means}: Gold-standard benchmark metric; high-quality production reward models achieve 70\%--85\% on held-out human comparisons.
    \item \textbf{Why It Is Important}: Direct measure of preference ranking fidelity before deploying to downstream policy optimization.
\end{itemize}

\subsection{Length Bias Pearson Correlation Monitor}
{\boldmath
\begin{equation}
\rho_{\text{len}} = \frac{\operatorname{Cov}(r_w - r_l, \; |y_w| - |y_l|)}{\sigma_{\Delta r} \sigma_{\Delta \text{len}}}
\end{equation}
}
\begin{itemize}
    \item \textbf{Formal Definition}: Normalized covariance between score difference and token count difference.
    \item \textbf{What It Means}: Diagnoses whether the reward model is evaluating quality or merely exploiting verbosity ($\rho > 0.4$ indicates severe length bias).
    \item \textbf{Why It Is Important}: Essential telemetry for preventing policy degeneration and verbose hallucination during RLHF.
\end{itemize}

\section{Rigorous Mathematical Analysis Checklist}

\begin{itemize}
    \item \textbf{Notation}: $x$ (prompt), $y_w$ (chosen completion), $y_l$ (rejected completion), $r_\theta \in \mathbb{R}$ (scalar reward), $m \ge 0$ (margin).
    \item \textbf{Epistemic Status}: Classical statistical model by Bradley and Terry (1952); adapted for modern LLMs by Christiano et al. (2017) and Ouyang et al. (2022).
    \item \textbf{Dependency Graph}: Dialogue Pair $(x, y_w, y_l) \to \text{Scalar Head} \to (r_w, r_l) \to \text{Sigmoid} \to \nabla_\theta \mathcal{L}_{\text{RM}} \to \text{RLHF Reward}$.
    \item \textbf{Dimensions}: Batch size $B$, reward scalar $r \in \mathbb{R}$.
    \item \textbf{Regimes}: RLHF alignment, Best-of-N response reranking, safety filtering, automated hallucination detection.
    \item \textbf{Invariants}: If $r_w - r_l \to +\infty$, then $P(y_w \succ y_l) \to 1.0$ and $\mathcal{L}_{\text{RM}} \to 0.0$.
    \item \textbf{Isomorphisms}: Bradley-Terry preference modeling is isomorphic to binary logistic regression on pairwise difference vectors.
\end{itemize}

\section{Theoretical Theorems and Formal Proofs}

\begin{theorem}[Convexity of the Bradley-Terry Objective with Respect to Latent Rewards]
The Bradley-Terry negative log-likelihood loss function $\mathcal{L}(\Delta r) = -\ln \sigma(\Delta r)$ is strictly convex with respect to the scalar reward margin $\Delta r = r_w - r_l$.
\end{theorem}

\begin{proof}
Let $z = \Delta r - m$. We write the loss function in terms of $z$:
{\boldmath
\begin{equation}
\ell(z) = -\ln \sigma(z) = -\ln \left( \frac{1}{1 + e^{-z}} \right) = \ln(1 + e^{-z})
\end{equation}
}
We evaluate the first derivative with respect to $z$:
{\boldmath
\begin{equation}
\ell'(z) = \frac{1}{1 + e^{-z}} \cdot (-e^{-z}) = -\frac{e^{-z}}{1 + e^{-z}} = -(1 - \sigma(z)) = \sigma(z) - 1
\end{equation}
}
We evaluate the second derivative with respect to $z$:
{\boldmath
\begin{equation}
\ell''(z) = \frac{d}{dz} (\sigma(z) - 1) = \sigma(z)(1 - \sigma(z))
\end{equation}
}
For all finite values of $z \in \mathbb{R}$, the logistic sigmoid function strictly satisfies $0 < \sigma(z) < 1$.
Therefore:
{\boldmath
\begin{equation}
\ell''(z) = \sigma(z)(1 - \sigma(z)) > 0, \quad \forall z \in \mathbb{R}
\end{equation}
}
Because the second derivative is strictly positive everywhere on $\mathbb{R}$, $\ell(z)$ is strictly convex.
Furthermore, the composition with the affine mapping $z = r_w - r_l - m$ preserves convexity, guaranteeing that the objective possesses no local minima with respect to the latent reward differences.
\end{proof}

\begin{theorem}[Translation Invariance of Bradley-Terry Preferences]
For any constant scalar $c \in \mathbb{R}$, shifting all reward evaluations uniformly $r'(x, y) = r(x, y) + c$ leaves the preference probability $P(y_w \succ y_l \mid x)$ and the training loss $\mathcal{L}_{\text{RM}}$ strictly invariant.
\end{theorem}

\begin{proof}
Let $r'(x, y) = r(x, y) + c$.
Evaluating the shifted preference difference:
{\boldmath
\begin{align}
\Delta r' &= r'(x, y_w) - r'(x, y_l) \\
&= (r(x, y_w) + c) - (r(x, y_l) + c) \\
&= r(x, y_w) - r(x, y_l) = \Delta r
\end{align}
}
Because the preference probability $P(y_w \succ y_l \mid x) = \sigma(\Delta r)$ and the loss $\mathcal{L}_{\text{RM}} = -\ln \sigma(\Delta r - m)$ depend solely on the scalar difference $\Delta r$, both quantities are identical under uniform shifts. This proves that Bradley-Terry models determine relative preference rankings rather than absolute scale anchors.
\end{proof}

\section{Foundational Architectural Questions and Epistemic Meta-Questions}

\subsection{Architectural Question 1: Reward Hacking and Over-Optimization}
\textbf{Question:} Why does optimizing a policy against a learned reward model lead to ``reward hacking'' at large optimization steps?\\
\textbf{Answer:} Goodhart's Law: ``When a measure becomes a target, it ceases to be a good measure.'' A neural reward model is an imperfect proxy for true human values with blind spots. Policy optimizers (like PPO) aggressively exploit these blind spots, generating bizarre formatting, repetitive phrases, or sycophantic rhetoric that maximizes the proxy score while degrading human-evaluated quality.

\textbf{Meta-Question:} How can reward hacking be detected in automated CI/CD pipelines?\\
\textbf{Answer:} CI/CD pipelines track three gold metrics: (1) KL divergence from the frozen base model; (2) Average response length drift; and (3) Periodic automated validation against a frozen panel of diverse judge models (e.g., ensemble of independent LLM-as-a-judge evaluators).

\subsection{Architectural Question 2: Margin-Based Reward Modeling}
\textbf{Question:} Why is a positive target margin $m > 0$ beneficial in Bradley-Terry training?\\
\textbf{Answer:} Standard cross-entropy treats any infinitesimally small margin $r_w > r_l$ as correct. Adding a target margin $m$ forces the network to establish a decisive, confident separation between preferred and rejected completions, improving generalization and calibration on subtle distinctions.

\textbf{Meta-Question:} Should reward models use generative reasoning heads before outputting a score?\\
\textbf{Answer:} Yes. Modern generative reward models (Critic / Reasoning RM) output a chain of critical reasoning explaining why response A is superior before emitting the final scalar reward, achieving substantially higher accuracy than simple regression heads.

\section{Algorithmic Implementation Specification}

\begin{algorithm}[H]
\caption{Bradley-Terry Preference Loss and Telemetry Evaluation}
\begin{algorithmic}[1]
\Require Chosen rewards $\mathbf{r}_w$, rejected rewards $\mathbf{r}_l$, lengths $\mathbf{L}_w, \mathbf{L}_l$, target margin $m$
\Ensure Loss $\mathcal{L}_{\text{RM}}$, pairwise accuracy $\operatorname{Acc}$, mean margin $\bar{\Delta r}$, length correlation $\rho_{\text{len}}$
\State Initialize $loss\_accum \gets 0.0, \quad correct \gets 0, \quad margin\_accum \gets 0.0$
\State Allocate $\boldsymbol{\Delta r} \in \mathbb{R}^B, \quad \boldsymbol{\Delta L} \in \mathbb{R}^B$
\For{$i = 0$ \textbf{to} $B-1$}
    \State $diff \gets \mathbf{r}_w[i] - \mathbf{r}_l[i]$
    \State $\boldsymbol{\Delta r}[i] \gets diff, \quad margin\_accum \gets margin\_accum + diff$
    \If{$diff > 0.0$}
        \State $correct \gets correct + 1$
    \EndIf
    \State $z \gets diff - m$
    \If{$z \ge 0.0$}
        \State $loss\_accum \gets loss\_accum + \ln(1 + \exp(-z))$
    \Else
        \State $loss\_accum \gets loss\_accum + (-z + \ln(1 + \exp(z)))$
    \EndIf
    \State $\boldsymbol{\Delta L}[i] \gets \mathbf{L}_w[i] - \mathbf{L}_l[i]$
\EndFor
\State $\mathcal{L}_{\text{RM}} \gets loss\_accum / B, \quad \operatorname{Acc} \gets correct / B, \quad \bar{\Delta r} \gets margin\_accum / B$
\State $\bar{\Delta L} \gets \frac{1}{B}\sum_{i=0}^{B-1} \boldsymbol{\Delta L}[i]$
\State $cov \gets \frac{1}{B}\sum_{i=0}^{B-1} (\boldsymbol{\Delta r}[i] - \bar{\Delta r})(\boldsymbol{\Delta L}[i] - \bar{\Delta L})$
\State $\sigma_r \gets \sqrt{\frac{1}{B}\sum (\boldsymbol{\Delta r}[i] - \bar{\Delta r})^2}, \quad \sigma_L \gets \sqrt{\frac{1}{B}\sum (\boldsymbol{\Delta L}[i] - \bar{\Delta L})^2}$
\State $\rho_{\text{len}} \gets cov / (\sigma_r \cdot \sigma_L + 10^{-9})$
\State \Return $(\mathcal{L}_{\text{RM}}, \operatorname{Acc}, \bar{\Delta r}, \rho_{\text{len}})$
\end{algorithmic}
\end{algorithm}

\section{Big-$\Theta$ Complexity Analysis}

\begin{itemize}
    \item \textbf{Bradley-Terry Loss Evaluation Time}: $\Theta(B)$ scalar softplus operations.
    \item \textbf{Pairwise Accuracy Time}: $\Theta(B)$ sign comparisons.
    \item \textbf{Length Correlation Analysis Time}: $\Theta(B)$ covariance and variance reductions.
    \item \textbf{Space Complexity}: $\Theta(B)$ buffer for difference metrics.
    \item \textbf{Parallel Reduction Depth}: $\mathcal{D} = \Theta(\log B)$ for empirical mean computation.
\end{itemize}

\section{Single Canonical Repository Reference}

The complete deterministic scalar implementation, verification suite, and unit test benchmarks are published at:
\begin{center}
\url{https://github.com/Chief-Strategist-J/llm-obs-infra}
\end{center}

\end{document}
